% Lösungen Bruchrechnung (zusammenbau v0.4, Heft)
\einheitenkopf{Bruchrechnung}

\erg{1}{a) $\frac{5}{6}$ \quad b) $\frac{7}{11}$ \quad c) $\frac{8}{14} = \frac{4}{7}$}

1a)

\bruchrechteck{5}{6}

\erg{2}{a) $\frac{25}{100} + \frac{9}{100} + \frac{4}{100} = \frac{38}{100} = \frac{19}{50}$ \quad b) $\frac{36}{100} + \frac{9}{100} + \frac{1}{100} = \frac{46}{100} = \frac{23}{50}$ \quad c) $\frac{8}{72} = \frac{1}{9}$ und $\frac{56}{504} = \frac{1}{9}$; $\frac{1}{9} + \frac{1}{9} + \frac{1}{9} = \frac{3}{9} = \frac{1}{3}$ \quad d) $\frac{11}{132} = \frac{1}{12}$ und $\frac{110}{1320} = \frac{1}{12}$; $\frac{1}{12} + \frac{1}{12} + \frac{1}{12} = \frac{3}{12} = \frac{1}{4}$ \quad e) $\frac{3}{8} + \frac{1}{8} = \frac{4}{8} = \frac{1}{2}$ \quad f) $\frac{27}{64} + \frac{27}{64} = \frac{54}{64} = \frac{27}{32}$ \quad g) $\frac{3}{12} + \frac{3}{12} = \frac{6}{12} = \frac{1}{2}$ \quad h) $\frac{7}{56} + \frac{7}{56} = \frac{14}{56} = \frac{1}{4}$ \quad i) $1 - \frac{4}{30} = \frac{26}{30} = \frac{13}{15}$ \quad j) $1 - \frac{6}{48} = \frac{42}{48} = \frac{7}{8}$ \quad k) $1 - \frac{3}{10} = \frac{7}{10}$ \quad l) $1 - \frac{2}{12} = \frac{10}{12} = \frac{5}{6}$}
\erg{3}{a) $4{,}36$ \quad b) $5{,}9$ \quad c) $5{,}60 + 1{,}23 = 6{,}83$}
\erg{4}{a) $2\,016{,}35\text{ €} - 1\,847{,}60\text{ €} = 168{,}75\text{ €}$ \quad b) $1\,004{,}15\text{ €} - 912{,}40\text{ €} = 91{,}75\text{ €}$}
\erg{5}{a) ja \quad b) $\frac{8}{9}$ \quad c) $\frac{6}{35}$}
\erg{6}{a) $\frac{7}{15} \cdot \frac{6}{14} \cdot \frac{5}{13} = \frac{210}{2730} = \frac{1}{13}$; $\frac{5}{15} \cdot \frac{4}{14} \cdot \frac{3}{13} = \frac{60}{2730} = \frac{2}{91}$; $210 : 60 = 3{,}5$, also dreieinhalbmal so wahrscheinlich, nicht doppelt. \quad b) $\frac{6}{10} \cdot \frac{5}{9} \cdot \frac{4}{8} = \frac{120}{720} = \frac{1}{6}$; $\frac{4}{10} \cdot \frac{3}{9} \cdot \frac{2}{8} = \frac{24}{720} = \frac{1}{30}$; $120 : 24 = 5$, also fünfmal so wahrscheinlich, nicht dreimal. \quad c) $\frac{9}{12} \cdot \frac{8}{11} = \frac{72}{132} = \frac{6}{11}$ \quad d) $\frac{7}{10} \cdot \frac{6}{9} = \frac{42}{90} = \frac{7}{15}$ \quad e) $\frac{24}{2184} = \frac{1}{91}$; das sind etwa $1{,}1\,\%$, also nicht kleiner als $1\,\%$. \quad f) $\frac{6}{3360} = \frac{1}{560}$; das sind etwa $0{,}18\,\%$, also kleiner als $1\,\%$. \quad g) $\frac{5}{6} \cdot \frac{1}{5} = \frac{5}{30} = \frac{1}{6}$ \quad h) $\frac{8}{10} \cdot \frac{2}{9} = \frac{16}{90} = \frac{8}{45}$}
\erg{7}{a) $37{,}2$ \quad b) $4{,}57$}
\erg{8}{a) $25\,000 : 100 \cdot 8 = 2\,000$ l; $2\,000 \cdot 1{,}624\text{ €} = 3\,248\text{ €}$ \quad b) $4\,000 : 100 \cdot 3 = 120$ l; $120 \cdot 1{,}785\text{ €} = 214{,}20\text{ €}$ \quad c) $30 \cdot 0{,}35\text{ ct} = 10{,}5\text{ ct}$ \quad d) $24 \cdot 0{,}26\text{ ct} = 6{,}24\text{ ct}$ \quad e) $0{,}2^2 = 0{,}04$ \quad f) $0{,}5^2 = 0{,}25$}
\erg{9}{a) die Malaufgabe einkreisen \quad b) $5 + 24 = 29$ \quad c) $5 \cdot 4 = 20$}
\erg{10}{a) $(4{,}5 + 1{,}5) : 4 = 6 : 4 = 1{,}5$ \quad b) $\left(\frac{4}{6} + \frac{1}{6}\right) : 5 = \frac{5}{6} : 5 = \frac{5}{30} = \frac{1}{6}$ \quad c) $(7 + 2{,}6) : 1{,}2 = 9{,}6 : 1{,}2 = 8$ \quad d) $(10 + 3{,}5) : 0{,}5 = 13{,}5 : 0{,}5 = 27$ \quad e) erster Term ja, zweiter Term nein (die Oma ist wie ein Kind gerechnet), dritter Term ja \quad f) erster Term nein (der Opa ist wie ein Kind gerechnet), zweiter Term ja, dritter Term ja}
